% element: 10-s002:e01
% element: 10-s002:e02
% element: 10-s002:e03
\slajdid{10-s002}
\begin{frame}[label=10-s002]{Model diode u logičkim kolima}
\setlength{\abovedisplayskip}{4pt}\setlength{\belowdisplayskip}{4pt}\setlength{\abovedisplayshortskip}{3pt}\setlength{\belowdisplayshortskip}{3pt}
\begin{columns}[T,onlytextwidth]\column{.47\textwidth}\centering
\begin{tikzpicture}[x=.85cm,y=.65cm,font=\sitneoznake,>=Latex]\ctikzset{bipoles/length=.65cm}\draw(-1,0)--(-.3,0);\draw(-.3,-.3)--(-.3,.3)--(.3,0)--cycle;\draw(.3,0)--(1,0);\draw(.3,-.3)--(.3,.3);\node[below]at(-.55,0){A};\node[below]at(.55,0){C};\draw[->](-1.5,.4)--(-1.05,.4)node[midway,above]{$I_D$};\draw[<->](-.75,.8)--(.65,.8)node[midway,above]{$+\ V_{AC}$};\draw(-.75,.6)--(-.75,1);\draw(.65,.6)--(.65,1);\node at(.12,.4){p n};\end{tikzpicture}
\par\medskip
A — anoda; C — katoda\par\bigskip
\Oznaka{Ne vodi}\medskip
\begin{tikzpicture}[x=.85cm,y=.65cm,font=\sitneoznake,>=Latex]\ctikzset{bipoles/length=.65cm}\draw(-1,0)node[ocirc]{}node[left]{A}--(-.25,0);\draw(.3,0)--(1,0)node[ocirc]{}node[right]{C};\draw[->](-.85,.3)--(-.4,.3)node[midway,above]{$I_D$};\end{tikzpicture}
\par
$V_{AC}<V_{\gamma D}$\column{.49\textwidth}\centering \Oznaka{Vodi}\medskip
\begin{tikzpicture}[x=.85cm,y=.65cm,font=\sitneoznake,>=Latex]\ctikzset{bipoles/length=.65cm}\draw(-1,0)node[ocirc]{}node[left]{A}--(-.08,0);\draw(.08,0)--(1,0)node[ocirc]{}node[right]{C};\draw(-.08,-.5)--(-.08,.5);\draw(.08,-.2)--(.08,.2);\node[above]at(0,.65){$V_D$};\node at(-.3,.35){$+$};\draw[->](-.85,.3)--(-.4,.3)node[midway,above]{$I_D$};\end{tikzpicture}
\par\medskip
Ako postoje uslovi za $V_{AC}>V_{\gamma D}$.\par\medskip
Ivica provođenja: $V_{AC}=V_{\gamma D}$, $I_D\approx0$.\end{columns}\bigskip
\Rezultat{Pretpostavka $\longrightarrow$ dokaz: ne vodi — proveri $V_{AC}<V_{\gamma D}$; vodi — proveri $I_D\geq0$.}
\end{frame}
